% GENERATED FILE. Edit canonical lessons, then run npm run generate:books.
% canonical-source-hash: fnv1a64:580d014ab338fae6
% canonical-lessons: GU-C208-sadak, GU-C208-khado, GU-C208-traphik, GU-C208-paidu, GU-C208-phatak

\chapter{Out on the highway}
\label{ch:gu-out-on-the-highway}
\hlchaptermodality{208}

\begin{chapteropening}
\textbf{By the end of this chapter:} \emph{I can describe the road we drove along.}

Complete the last lesson of chapter 208: I can describe the road we drove along.
\end{chapteropening}

\section[\texorpdfstring{\gu{સડક} (saḍak)}{સડક (saḍak)}]{\gu{સડક} (saḍak) — main road, paved road}

\label{lesson:GU-C208-sadak}



\begin{quote}
Before the new one: say the Gujarati for petrol, then the Gujarati for to hire.
\end{quote}

\subsection*{You'll want to know: \gu{સડક}}
\textbf{\gu{સડક}} — \emph{saḍak} — \textquotedblleft{}main road, paved road\textquotedblright{}.

\begin{quote}
\textbf{\gu{સડક} \gu{પહોળી} \gu{હતી}.} — \emph{saḍak pahoḷī hatī} — \textquotedblleft{}The road was wide.\textquotedblright{}
\end{quote}

\begin{cousinweb}
\textbf{\gu{સડક}} (\emph{saḍak}) is a tadbhava. CDIAL lists it, glossed \textquotedblleft{}road\textquotedblright{}, under Sanskrit \emph{sṛtí}, \textquotedblleft{}road, path\textquotedblright{}, extended with the suffix \emph{-kk-} (CDIAL 13577).
\end{cousinweb}

\begin{grammarlens}[title={What you've built}]
One more word for everyday life at A2.
\end{grammarlens}

\subsection*{Your turn}
\begin{itemize}
\raggedright
  \item \emph{Say it:} \emph{saḍak}
  \item \emph{Say it:} \emph{saḍak}, once more
  \item \emph{Say it:} \emph{saḍak pahoḷī hatī} — The road was wide.
  \item \emph{Recall:} say the Gujarati for petrol, then the Gujarati for to hire, then say \emph{saḍak} again
\end{itemize}

\begin{tcolorbox}[breakable,colback=teal!4,colframe=teal!35!black,title={Before you move on}]
What does \gu{સડક} mean? (\textquotedblleft{}Main road, paved road\textquotedblright{}.) Say it once more.
\end{tcolorbox}
\section[\texorpdfstring{\gu{ખાડો} (khāḍo)}{ખાડો (khāḍo)}]{\gu{ખાડો} (khāḍo) — pothole, pit}

\label{lesson:GU-C208-khado}



\begin{quote}
Before the new one: say the Gujarati for to hire, then the Gujarati for main road.
\end{quote}

\subsection*{You'll want to know: \gu{ખાડો}}
\textbf{\gu{ખાડો}} — \emph{khāḍo} — \textquotedblleft{}pothole, pit\textquotedblright{}.

\begin{quote}
\textbf{\gu{રસ્તામાં} \gu{મોટો} \gu{ખાડો} \gu{હતો}.} — \emph{rastāmā̃ moṭo khāḍo hato} — \textquotedblleft{}There was a big pothole in the road.\textquotedblright{}
\end{quote}

\begin{cousinweb}
\textbf{\gu{ખાડો}} (\emph{khāḍo}) is a tadbhava. CDIAL lists it, glossed \textquotedblleft{}hole\textquotedblright{}, under \emph{*khāḍa}, \textquotedblleft{}a hollow\textquotedblright{}, a reconstructed form that it compares with \emph{*khaḍḍa}, \textquotedblleft{}hole, pit\textquotedblright{} (CDIAL 3860).
\end{cousinweb}

\begin{grammarlens}[title={What you've built}]
One more word for everyday life at A2.
\end{grammarlens}

\subsection*{Your turn}
\begin{itemize}
\raggedright
  \item \emph{Say it:} \emph{khāḍo}
  \item \emph{Say it:} \emph{khāḍo}, once more
  \item \emph{Say it:} \emph{rastāmā̃ moṭo khāḍo hato} — There was a big pothole in the road.
  \item \emph{Recall:} say the Gujarati for to hire, then the Gujarati for main road, then say \emph{khāḍo} again
\end{itemize}

\begin{tcolorbox}[breakable,colback=teal!4,colframe=teal!35!black,title={Before you move on}]
What does \gu{ખાડો} mean? (\textquotedblleft{}Pothole, pit\textquotedblright{}.) Say it once more.
\end{tcolorbox}
\section[\texorpdfstring{\gu{ટ્રાફિક} (ṭrāphik)}{ટ્રાફિક (ṭrāphik)}]{\gu{ટ્રાફિક} (ṭrāphik) — traffic}

\label{lesson:GU-C208-traphik}



\begin{quote}
Before the new one: say the Gujarati for main road, then the Gujarati for pothole.
\end{quote}

\subsection*{You'll want to know: \gu{ટ્રાફિક}}
\textbf{\gu{ટ્રાફિક}} — \emph{ṭrāphik} — \textquotedblleft{}traffic\textquotedblright{}.

\begin{quote}
\textbf{\gu{શહેરમાં} \gu{ઘણો} \gu{ટ્રાફિક} \gu{હતો}.} — \emph{śahermā̃ ghaṇo ṭrāphik hato} — \textquotedblleft{}There was a lot of traffic in the city.\textquotedblright{}
\end{quote}

\begin{cousinweb}
\textbf{\gu{ટ્રાફિક}} (\emph{ṭrāphik}) is a loan from English \emph{traffic}.
\end{cousinweb}

\begin{grammarlens}[title={What you've built}]
One more word for everyday life at A2.
\end{grammarlens}

\subsection*{Your turn}
\begin{itemize}
\raggedright
  \item \emph{Say it:} \emph{ṭrāphik}
  \item \emph{Say it:} \emph{ṭrāphik}, once more
  \item \emph{Say it:} \emph{śahermā̃ ghaṇo ṭrāphik hato} — There was a lot of traffic in the city.
  \item \emph{Recall:} say the Gujarati for main road, then the Gujarati for pothole, then say \emph{ṭrāphik} again
\end{itemize}

\begin{tcolorbox}[breakable,colback=teal!4,colframe=teal!35!black,title={Before you move on}]
What does \gu{ટ્રાફિક} mean? (\textquotedblleft{}Traffic\textquotedblright{}.) Say it once more.
\end{tcolorbox}
\section[\texorpdfstring{\gu{પૈડું} (paiḍũ)}{પૈડું (paiḍũ)}]{\gu{પૈડું} (paiḍũ) — wheel}

\label{lesson:GU-C208-paidu}



\begin{quote}
Before the new one: say the Gujarati for pothole, then the Gujarati for traffic.
\end{quote}

\subsection*{You'll want to know: \gu{પૈડું}}
\textbf{\gu{પૈડું}} — \emph{paiḍũ} — \textquotedblleft{}wheel\textquotedblright{}.

\begin{quote}
\textbf{\gu{ગાડીનું} \gu{એક} \gu{પૈડું} \gu{ખાડામાં} \gu{ગયું}.} — \emph{gāḍīnũ ek paiḍũ khāḍāmā̃ gayũ} — \textquotedblleft{}One wheel of the car went into the pothole.\textquotedblright{}
\end{quote}

\begin{cousinweb}
\textbf{\gu{પૈડું}} (\emph{paiḍũ}) is a tadbhava. CDIAL lists it, glossed \textquotedblleft{}felly\textquotedblright{}, the rim of a wheel, under Sanskrit \emph{pradhí}, \textquotedblleft{}felly\textquotedblright{}, extended with the suffix \emph{-ḍa-} (CDIAL 8672).
\end{cousinweb}

\begin{grammarlens}[title={What you've built}]
One more word for everyday life at A2.
\end{grammarlens}

\subsection*{Your turn}
\begin{itemize}
\raggedright
  \item \emph{Say it:} \emph{paiḍũ}
  \item \emph{Say it:} \emph{paiḍũ}, once more
  \item \emph{Say it:} \emph{gāḍīnũ ek paiḍũ khāḍāmā̃ gayũ} — One wheel of the car went into the pothole.
  \item \emph{Recall:} say the Gujarati for pothole, then the Gujarati for traffic, then say \emph{paiḍũ} again
\end{itemize}

\begin{tcolorbox}[breakable,colback=teal!4,colframe=teal!35!black,title={Before you move on}]
What does \gu{પૈડું} mean? (\textquotedblleft{}Wheel\textquotedblright{}.) Say it once more.
\end{tcolorbox}
\section[\texorpdfstring{\gu{ફાટક} (phāṭak)}{ફાટક (phāṭak)}]{\gu{ફાટક} (phāṭak) — level crossing, gate}

\label{lesson:GU-C208-phatak}



\begin{quote}
Before the new one: say the Gujarati for traffic, then the Gujarati for wheel.
\end{quote}

\subsection*{You'll want to know: \gu{ફાટક}}
\textbf{\gu{ફાટક}} — \emph{phāṭak} — \textquotedblleft{}level crossing, gate\textquotedblright{}.

\begin{quote}
\textbf{\gu{ફાટક} \gu{બંધ} \gu{હતું}.} — \emph{phāṭak bandh hatũ} — \textquotedblleft{}The level-crossing gate was closed.\textquotedblright{}
\end{quote}

A neuter noun.

\begin{cousinweb}
\textbf{\gu{ફાટક}} (\emph{phāṭak}) is usually traced to \emph{*phāṭṭakka}, \textquotedblleft{}gate, door\textquotedblright{}, a reconstructed form whose own source CDIAL leaves in doubt. CDIAL lists Hindi \emph{phāṭak}, \textquotedblleft{}gate\textquotedblright{}, but for Gujarati only \emph{phaṭkiyũ bārṇũ}, \textquotedblleft{}a single shutter to close a door\textquotedblright{} (CDIAL 9069).
\end{cousinweb}

\begin{grammarlens}[title={What you've built}]
One more word for everyday life at A2.
\end{grammarlens}

\subsection*{Your turn}
\begin{itemize}
\raggedright
  \item \emph{Say it:} \emph{phāṭak}
  \item \emph{Say it:} \emph{phāṭak}, once more
  \item \emph{Say it:} \emph{phāṭak bandh hatũ} — The level-crossing gate was closed.
  \item \emph{Recall:} say the Gujarati for traffic, then the Gujarati for wheel, then say \emph{phāṭak} again
\end{itemize}

\begin{tcolorbox}[breakable,colback=teal!4,colframe=teal!35!black,title={Before you move on}]
What does \gu{ફાટક} mean? (\textquotedblleft{}Level crossing, gate\textquotedblright{}.) Say it once more.
\end{tcolorbox}
